% !TeX program = xelatex
% Compile:  xelatex main.tex   (run twice)   — requires the xepersian package
% Fonts: B Nazanin (standard in Iran). Alternatives: XB Zar, XB Niloofar
\documentclass[12pt,a4paper]{article}
\usepackage[margin=2.3cm]{geometry}
\usepackage{amsmath,amssymb,graphicx,float,booktabs,longtable}
\usepackage{xepersian}
\settextfont{B Nazanin}
\setlatintextfont[BoldFont={Times New Roman}, ItalicFont={Times New Roman}, BoldItalicFont={Times New Roman}]{Times New Roman}
\begin{document}

\begin{center}
{\Large \textbf{ارائه یک روش بهبود تصویر کم‌نور مبتنی بر پیشین (\lr{Prior)} مدل‌های انتشار (\lr{Diffusion)} برای شرایط نوری نامتوازن گرگ‌ومیش با کاربرد در بینایی ماشین}}\\[6pt]
{\normalsize \lr{A Diffusion-Prior Zero-Shot Method for Low-Light Image Enhancement under Imbalanced Twilight Conditions with Application to Machine Vision}}
\end{center}
\noindent\textbf{کلیدواژه‌ها:} بهبود تصویر کم‌نور، مدل‌های انتشار، یادگیری بدون نظارت، گرگ‌ومیش، بینایی ماشین\\
\noindent\textbf{\lr{Keywords:}} \lr{Low-Light Image Enhancement, Diffusion Models, Zero-Shot Learning, Twilight, Machine Vision}


\section{شرح مسئله (اهداف، سابقه و ضرورت تحقیق)}
\textbf{۱-۱. اهمیت کاربردی بهبود تصویر کم‌نور در بینایی ماشین}\par
درصد قابل‌توجهی از کاربردهای حیاتی بینایی ماشین، از رانندگی خودران و نظارت تصویری تا عکاسی محاسباتی و پایش محیطی، ناگزیر باید در شرایط نوری نامطلوب، و به‌ویژه در غروب، شب و فضاهای سرپوشیده کم‌نور، عمل کنند [\lr{1}]. در چنین شرایطی، محدودیت‌های فیزیکی حسگرهای تصویری (اندازه پیکسل، ظرفیت چاه الکترون و عملکرد سنسور در \lr{ISO}های بالا) موجب می‌شود تصویر ثبت‌شده با شدت روشنایی ناکافی، نویز شدید، اغتشاش رنگ و افت کنتراست همراه باشد؛ به‌طوری‌که هم کیفیت ادراکی تصویر برای ناظر انسانی مخدوش می‌شود و هم عملکرد مدل‌های سطح بالا مانند تشخیص شیء، قطعه‌بندی معنایی و تخمین عمق که عمدتاً روی تصاویر نور طبیعی آموزش دیده‌اند، به‌طور چشمگیری افت می‌کند [\lr{2}, \lr{3}]. مطالعه مرجع [\lr{2}] نشان داده است که دقت تشخیص چهره در تصاویر کم‌نور می‌تواند به‌مراتب کمتر از حالت نور طبیعی باشد، و همین شکاف، بهبود تصویر کم‌نور (\lr{Low-Light} \lr{Image} \lr{Enhancement}: \lr{LLIE)} را به یک پیش‌نیاز عملی برای زنجیره کامل پردازش بصری تبدیل کرده است.\par
اهمیت این مسئله در حوزه رانندگی هوشمند مضاعف است. دیتاست‌های \lr{adverse-condition} مانند \lr{ACDC} [\lr{4}] و \lr{DarkDriving} [\lr{5}] نشان می‌دهند که صحنه‌های شب و گرگ‌ومیش، سهم عمده‌ای از شرایط عملیاتی واقعی خودروهای خودران را تشکیل می‌دهند و خطای ادراک در این رژیم‌ها مستقیماً به ریسک ایمنی ترجمه می‌شود. از سوی دیگر، در نظارت تصویری شهری و صنعتی، اکثر رخدادهای امنیتی در ساعات تاریکی رخ می‌دهند و کیفیت پایین تصاویر، بازسازی رخداد را دشوار می‌سازد [\lr{1}]. در عکاسی محاسباتی نیز خط لوله دوربین (\lr{ISP)} در نوردهی کوتاه و \lr{ISO} بالا با اغتشاش‌های خام مواجه است و رویکرد تعویض مدل مولد برای \lr{ISP} دوربین، هم‌اکنون در خط مقدم صنعت قرار دارد [\lr{6}].\par
از منظر روش‌شناسی، دو راهکار موازی برای مواجهه با نوریافتگی ناکافی وجود دارد: (الف) مداخله در زمان ثبت، مانند افزایش زمان نوردهی یا \lr{ISO}، که به‌ترتیب به لرزش و نویز منجر می‌شود، و (ب) پردازش پس از ثبت با الگوریتم‌های \lr{LLIE.} راهکار دوم به دلیل عدم نیاز به سخت‌افزار ویژه و قابلیت افزودن به هر سامانه موجود، در پژوهش‌های اخیر کانون توجه بوده است. مجموعه مقالات مروری اخیر، رشد انفجاری این حوزه را در پنج سال گذشته نشان می‌دهند [\lr{7}, \lr{1}]. شکل ۱ نمونه‌هایی از پیچیدگی‌های واقعی تصاویر کم‌نور دنیا (شامل نوردهی هم‌زمان کم و زیاد در یک قاب) را نمایش می‌دهد و شکل ۲ سیر تحول روش‌های این حوزه را در یک نگاه تاریخی ترسیم می‌کند.\par
\begin{figure}[H]\centering
\includegraphics[width=0.92\linewidth]{figures/fig-01.png}\\[4pt]
\end{figure}
شکل ۱: نمونه تصاویر دنیای واقعی با نوردهی مخلوط از دیتاست \lr{SICE} (برگرفته از \lr{Fig.1} در مرجع [\lr{1}])\par
\begin{figure}[H]\centering
\includegraphics[width=0.92\linewidth]{figures/fig-02.png}\\[4pt]
\end{figure}
شکل ۲: تایم‌لاین روش‌های نماینده بهبود تصویر کم‌نور (برگرفته از \lr{Fig.2} در مرجع [\lr{1}])\par
مسئله‌ای که در ادامه خواهیم دید، این است که قریب به اتفاق روش‌ها و دیتاست‌های این حوزه بر «شب مطلق» متمرکز بوده‌اند و رژیم نوری واسط (به‌ویژه گرگ‌ومیش که هم‌زمان با روشنایی نسبی آسمان، سایه‌های عمیق و منابع نوری مصنوعی روبه‌روست) عملاً در بنچمارک‌ها غایب است. این غیاب، شکاف پژوهشی اصلی رساله حاضر را شکل می‌دهد (زیربخش ۱-۵).\par
\textbf{۱-۲. خط روش‌های نظارت‌شده: از تجزیه \lr{Retinex} تا ترنسفورمر و اشباع بنچمارک}\par
\textbf{۱-۲-۱. پیشینه کلاسیک: از تئوری \lr{Retinex} تا روش‌های بهینه‌سازی‌محور}\par
ریشه‌های مسئله به دهه ۱۹۷۰ بازمی‌گردد، زمانی که \lr{Land} و \lr{McCann} با طرح تئوری \lr{Retinex}، بازنمایی فیزیکی از ثبات رنگی بینایی انسان ارائه کردند: تصویر مشاهده‌شده حاصل ضرب مؤلفه بازتاب (\lr{Reflectance)}، ویژگی ذاتی صحنه، و مؤلفه روشنایی (\lr{Illumination)} است [\lr{8}]. این صورت‌بندی ساده، به‌رغم عدم یکتایی تجزیه (مسئله‌ای ذاتاً بدتعریف یا \lr{ill-posed)}، به چارچوب مفهومی غالب حوزه تبدیل شد و قریب به پنج دهه بعد، همچنان هسته بسیاری از معماری‌های امروزی، از \lr{RetinexNet} تا \lr{Retinexformer} و \lr{ReCo-Diff}، قابل شناسایی است.\par
در دو دهه بعد، تمرکز پژوهش بر روش‌های پردازشی مستقیم و بدون مدل یادگیری بود. تراز هیستوگرام تطبیقی و واریانت محبوب آن \lr{CLAHE} [\lr{9}] با بازتوزیع محلی هیستوگرام، کنتراست نواحی تاریک را افزایش می‌داد اما بهای آن، تقویت نویز و اغتشاش رنگ در نواحی ناهموار بود. در خط \lr{Retinex}، \lr{Jobson} و همکاران با \lr{SSR} و سپس نسخه چندمقیاسی \lr{MSRCR} [\lr{10}]، روشنایی را با فیلترهای گوسی چندمقیاس تخمین زدند و بازتاب رنگی را بازیابی کردند؛ رویکردی که کیفیت رنگی بهتری می‌داد اما به انتخاب دستی پارامترهای گوسی حساس بود.\par
دهه ۲۰۱۰ نقطه بلوغ روش‌های بهینه‌سازی‌محور بود. \lr{Dong} و همکاران [\lr{11}] با قیاس هوشمندانه میان کم‌نوری و مه، مدل برداشتن غبار را وارونه روی تصاویر تاریک اعمال کردند؛ ایده‌ای که به‌تنهایی نشان داد انتقال دانش میان مسائل \lr{low-level} می‌تواند راه‌گشا باشد. سپس مدل‌های واریانس‌پذیر وزن‌دار وارد شدند: \lr{SRIE} [\lr{12}] بازتاب و روشنایی را در قالب یک مسئله کمینه‌سازی وزن‌دار به‌طور هم‌زمان تخمین زد و \lr{LIME} [\lr{13}] با ساختارپذیرسازی قید روشنایی، به تخمین نگاشت روشنایی موفق‌تری رسید. وجه مشترک تمام این خانواده، قیدها و پارامترهای دست‌ساز بود: تفسیرپذیری بالایی داشتند، اما حل‌کننده‌های تکراریِ کند، حساسیت به تنظیم پارامتر و، مهم‌تر از همه، ناتوانی در تعمیم به صحنه‌های متنوع، سقف کیفیت‌شان را تعیین می‌کرد. جدول ۱ این مسیر تاریخی را با محدودیت اصلی هر دوره خلاصه می‌کند.\par
جدول ۱: سیر تحول رویکردهای کلاسیک بهبود تصویر کم‌نور و محدودیت اصلی هر دوره (تألیف پژوهش حاضر بر پایه مراجع ذکرشده)\par
\begin{table}[H]\centering
{\small
\begin{tabular}{|c|c|c|c|}
\hline
دوره & رویکرد & نمایندگان & محدودیت اصلی \\ \hline
۱۹۷۱ تا ۱۹۷۷ & تئوری \lr{Retinex} & \lr{Land} و \lr{McCann} [\lr{8}] & چارچوب نظری؛ حل تحلیلی به‌دلیل عدم یکتایی تجزیه ممکن نیست \\ \hline
دهه ۱۹۸۰ & تراز هیستوگرام تطبیقی & \lr{AHE/CLAHE} [\lr{9}] & تقویت نویز و اغتشاش رنگ در نواحی ناهموار \\ \hline
۱۹۹۷ & \lr{Retinex} چندمقیاس & \lr{SSR/MSRCR} [\lr{10}] & حساسیت به پارامترهای گوسی و رنگ مصنوعی \\ \hline
۲۰۱۱ & قیاس با برداشتن غبار & \lr{Dong} \lr{et} \lr{al.} [\lr{11}] & فرض مدل پراکندگی در صحنه کم‌نور واقعی برقرار نیست \\ \hline
۲۰۱۶ & مدل واریانس‌پذیر وزن‌دار & \lr{SRIE} [\lr{12}] & حل‌کننده تکراری کند؛ تنظیم دستی پارامترها \\ \hline
۲۰۱۷ & تخمین نگاشت روشنایی & \lr{LIME} [\lr{13}] & وابستگی به \lr{prior} ساختاری دست‌ساز \\ \hline
۲۰۱۷ & نخستین یادگیری عمیق & \lr{LLNet} [\lr{14}] & ظرفیت محدود \lr{autoencoder} و آموزش ناپایدار \\ \hline
\end{tabular}}
\end{table}
نقطه عطف بعدی، ورود یادگیری عمیق بود: \lr{LLNet} [\lr{14}] در ۲۰۱۷ با یک \lr{autoencoder} انباشته نشان داد مدل می‌تواند هم‌زمان تقویت و حذف نویز را از داده بیاموزد، هرچند ظرفیت محدود و آموزش ناپایدار، آن را به یک اثبات مفهومی بیش تبدیل نمی‌کرد. همین اثبات مفهومی بود که مقدمه انفجار ۲۰۱۸ به بعد شد؛ مرحله‌ای که در زیربخش بعدی به آن می‌پردازیم.\par
شکل ۳ این دسته‌بندی را در یک نگاه ترسیم می‌کند.\par
\begin{figure}[H]\centering
\includegraphics[width=0.92\linewidth]{figures/fig-03.png}\\[4pt]
\end{figure}
شکل ۳: رده‌بندی سلسله‌مراتبی استراتژی‌های یادگیری برای \lr{LLIE} (برگرفته از \lr{Fig.3} در مرجع [\lr{1}])\par
\textbf{۱-۲-۲. نسل اول شبکه‌های نظارت‌شده مبتنی بر \lr{Retinex}}\par
\lr{RetinexNet} نخستین چارچوب یادگیری‌محور برای تجزیه داده‌محور \lr{Retinex} بود؛ این روش با معرفی دیتاست \lr{LOL} (۵۰۰ جفت تصویر واقعی کم‌نور/عادی، ۴۸۵ برای آموزش و ۱۵ برای ارزیابی)، دو شبکه \lr{Decom-Net} و \lr{Enhance-Net} را سرتاسری آموزش داد [\lr{15}]. هم‌زمان، \lr{SID} نشان داد پردازش مستقیم داده خام (\lr{RAW)} سنسور با شبکه‌های یادگیری عمیق، کیفیت به‌مراتب بهتری نسبت به پردازش تصویر \lr{sRGB} تولید می‌کند، هرچند نیازمند جفت‌های نوردهی کوتاه/بلند گران‌قیمت است [\lr{16}]. \lr{KinD} با تفکیک آشکار نوردهی/حذف نویز و معماری چندمرحله‌ای، نتایج \lr{RetinexNet} را بهبود داد [\lr{17}] و نسخه ژورنالی آن (\lr{KinD++)} کیفیت بازسازی را باز هم بالاتر برد [\lr{18}]. روش‌های دیگری مانند \lr{DeepUPE} (تخمین مستقیم نگاشت روشنایی) [\lr{19}] و \lr{DRBN} (رویکرد نیمه‌نظارت‌شده مبتنی بر بازنمایی نوار) [\lr{20}] نیز به این خط افزوده شدند. در ادامه، \lr{RUAS} با بازکردن (\lr{Unrolling)} بهینه‌سازی \lr{Retinex} به حلقه‌های شبکه‌ای و جست‌وجوی معماری همکار، و \lr{URetinex-Net} با چارچوب \lr{unfold} عمیق، پایداری تجزیه را ارتقا دادند [\lr{21}, \lr{22}]؛ \lr{SNR-Net} نیز با وزن‌دهی آگاه از نسبت سیگنال‌به‌نویز، نویز تقویت‌شده در فرایند روشن‌سازی را مدیریت کرد [\lr{23}].\par
\begin{figure}[H]\centering
\includegraphics[width=0.92\linewidth]{figures/fig-04.png}\\[4pt]
\end{figure}
شکل ۴: معماری \lr{Retinex-Net} (برگرفته از \lr{Fig.1} در مرجع [\lr{15}])\par
\begin{figure}[H]\centering
\includegraphics[width=0.92\linewidth]{figures/fig-05.png}\\[4pt]
\end{figure}
شکل ۵: نمونه جفت‌های دیتاست \lr{LOL} (برگرفته از \lr{Fig.3} در مرجع [\lr{15}])\par
\textbf{۱-۲-۳. مدل‌های مولد و ترنسفورمر: قله خط نظارت‌شده}\par
\lr{LLFlow} با استدلالی مهم وارد میدان شد: نگاشت کم‌نور به عادی ذاتاً یک‌به‌چند است و تابع هزینه پیکسل‌محور \lr{L1} نمی‌تواند توزیع شرطی تصاویر با نوردهی عادی را مدل کند؛ این روش با \lr{conditional} \lr{normalizing} \lr{flow} و توزیع احتمال کامل، جهشی حدود ۴ دسی‌بل در \lr{PSNR} روی \lr{LOL} نسبت به روش‌های پیشین ایجاد کرد [\lr{24}].\par
\begin{figure}[H]\centering
\includegraphics[width=0.92\linewidth]{figures/fig-06.png}\\[4pt]
\end{figure}
شکل ۶: معماری \lr{LLFlow} مبتنی بر جریان نرمال‌کننده (برگرفته از \lr{Fig.2} در مرجع [\lr{24}])\par
این بینش (مدل‌سازی توزیع به‌جای نگاشت قطعی) از نظر مفهومی، پیش‌درآمد موج \lr{diffusion} است. سپس \lr{Retinexformer} با قاب یک‌مرحله‌ای مبتنی بر \lr{Retinex} (\lr{ORF)} و ترنسفورمر هدایت‌شده با روشنایی (\lr{IGT)}، هم خطای چندمرحله‌ای آموزش را حذف کرد و هم در سیزده بنچمارک، از جمله \lr{PSNR} حدود ۲۵/۲ در \lr{LOL-v1} و ۲۵/۷ در \lr{LOL-v2-synthetic}، وضعیت \lr{SOTA} را تثبیت نمود [\lr{25}].\par
در ادامه، \lr{HVI-CIDNet} با معرفی فضای رنگ جدید \lr{HVI} و جداسازی آگاهانه روشنایی/رنگ، اغتشاش‌های رنگی ناشی از روشن‌سازی را هدف گرفت و جایگاه این خط را تا \lr{CVPR} \lr{2025} تثبیت کرد [\lr{26}].\par
جدول ۲ پیشرفت کمّی این خط را روی سه نسخه بنچمارک \lr{LOL} خلاصه می‌کند؛ صعود نزدیک به ۸/۵ دسی‌بل \lr{PSNR} از \lr{RetinexNet} تا \lr{Retinexformer}، هم‌زمان با نشستن مدل‌ها روی توزیع باریک همین بنچمارک‌ها، تصویر روشنی از «اشباع مذکور» ترسیم می‌کند.\par
جدول ۲: مقایسه کمی روش‌های نظارت‌شده روی بنچمارک‌های \lr{LOL}؛ \lr{PSNR} (\lr{dB)} و \lr{SSIM} (اعداد گزارش‌شده در [\lr{25}])\par
\begin{table}[H]\centering
{\small
\begin{tabular}{|c|c|c|c|}
\hline
روش & \lr{LOL-v1} (\lr{PSNR/SSIM)} & \lr{LOL-v2-real} (\lr{PSNR/SSIM)} & \lr{LOL-v2-syn} (\lr{PSNR/SSIM)} \\ \hline
\lr{RetinexNet} & \lr{16.77/0.560} & \lr{15.47/0.567} & \lr{17.13/0.798} \\ \hline
\lr{DeepUPE} & \lr{14.38/0.446} & \lr{13.27/0.452} & \lr{15.08/0.623} \\ \hline
\lr{EnlightenGAN} & \lr{17.48/0.650} & \lr{18.23/0.617} & \lr{16.57/0.734} \\ \hline
\lr{RUAS} & \lr{18.23/0.720} & \lr{18.37/0.723} & \lr{16.55/0.652} \\ \hline
\lr{DRBN} & \lr{20.13/0.830} & \lr{20.29/0.831} & \lr{23.22/0.927} \\ \hline
\lr{KinD} & \lr{20.86/0.790} & \lr{14.74/0.641} & \lr{13.29/0.578} \\ \hline
\lr{Restormer} & \lr{22.43/0.823} & \lr{19.94/0.827} & \lr{21.41/0.830} \\ \hline
\lr{MIRNet} & \lr{24.14/0.830} & \lr{20.02/0.820} & \lr{21.94/0.876} \\ \hline
\lr{SNR-Net} & \lr{24.61/0.842} & \lr{21.48/0.849} & \lr{24.14/0.928} \\ \hline
\lr{Retinexformer} & \lr{25.16/0.845} & \lr{22.80/0.840} & \lr{25.67/0.930} \\ \hline
\end{tabular}}
\end{table}
\begin{figure}[H]\centering
\includegraphics[width=0.92\linewidth]{figures/fig-07.png}\\[4pt]
\end{figure}
شکل ۷: معماری \lr{Retinexformer} (برگرفته از \lr{Fig.2} در مرجع [\lr{25}])\par
\textbf{۱-۲-۴. اشباع خط نظارت‌شده و محدودیت‌های ساختاری آن}\par
با وجود اعداد چشمگیر، سه محدودیت ساختاری این خانواده قابل‌انکار نیست:\par
۱. وابستگی به جفت‌داده گران‌قیمت. کل این خانواده روی جفت‌های واقعی (\lr{LOL} با ۴۸۵ نمونه آموزش، \lr{LOL-v2} با ۶۸۹ نمونه [\lr{27}]) یا سنتزی آموزش می‌بینند؛ داده‌جمعی واقعی در رژیم‌های نوری خاص بسیار پرهزینه است و اندازه کوچک دیتاست، عملاً سقف تعمیم را تعیین می‌کند.\par
۲. افت تعمیم به رژیم‌های نوری دیده‌نشده. آموزش روی توزیع «شب مطلق» \lr{LOL} و ارزیابی روی همان توزیع، عملکرد در رژیم‌های واسط (مثل نوردهی نامتوازن گرگ‌ومیش) را تضمین نمی‌کند؛ پدیده‌ای که مقالات مروری اخیر نیز به‌صراحت گزارش کرده‌اند [\lr{1}].\par
۳. هزینه محاسباتی مدل‌های قله. روش‌هایی مانند \lr{Retinexformer} و \lr{LLFlow} برای استنتاج بلادرنگ روی دستگاه‌های محدود مناسب نیستند [\lr{25}].\par
جمع‌بندی این زیربخش برای روایت رساله: خط نظارت‌شده از نظر «حداکثر \lr{PSNR} روی دیتاست خودش» به اشباع رسیده و بهبودهای حاشیه‌ای، دیگر موضوع پژوهشی معناداری نیستند؛ مسئله واقعی، تعمیم‌پذیری به رژیم‌های نوری واقعی و دیده‌نشده بدون اتکا به جفت‌داده است؛ مسیری که خط \lr{unsupervised/zero-shot} (زیربخش ۱-۳) و سپس موج \lr{diffusion-prior} (زیربخش ۱-۴) باز کرده‌اند.\par
\textbf{۱-۳. خط روش‌های بدون جفت‌داده: \lr{unsupervised} و \lr{zero-shot}}\par
\textbf{۱-۳-۱. انگیزه: رهایی از قید جفت‌داده}\par
محدودیت‌های ساختاری خط نظارت‌شده (زیربخش ۱-۲-۴) پرسشی طبیعی پیش می‌گذارد: آیا می‌توان بدون داده جفتی (که جمع‌آوری آن هم پرهزینه است و هم توزیع آن به رژیم نوری خاصی قفل می‌شود) به بهبود معنادار دست یافت؟ پاسخ مثبت این پرسش، دو زیرخانواده را پدید آورد: روش‌های بدون‌جفت (\lr{unpaired)} مانند \lr{EnlightenGAN} که با مولد حاوی توجه و رقیب بدون جفت، بهبود طبیعی‌تری نسبت به روش‌های نظارت‌شده اولیه ارائه کرد [\lr{28}]) و روش‌های \lr{zero-reference/zero-shot.} این دو برچسب اگرچه در ادبیات گاهی به‌ج هم به کار می‌روند، یکسان نیستند و در این رساله تفکیک می‌شوند: روش‌های \lr{zero-reference} مانند \lr{Zero-DCE} هنوز شبکه‌ای را (صرفاً با قیدهای طراحی‌شده روی تصویر و بدون هیچ داده مرجع) آموزش می‌دهند، در حالی که روش‌های \lr{zero-shot} در زمان آزمون و بدون هرگونه آموزش، با بهره‌گیری از یک مدل پیش‌آموزش‌دیده با ورودی مواجه می‌شوند [\lr{29}].\par
\textbf{۱-۳-۲. روش‌های شاخص \lr{zero-reference}}\par
\lr{Zero-DCE} مسئله را به تخمین منحنی‌های تقویتی تصویر-ویژه تبدیل کرد؛ شبکه‌ای فوق سبک (حدود ۸۰ هزار پارامتر) که با قیدهای بدون مرجع (فضایی-یکنواختی، نوردهی و حفظ رنگ) آموزش می‌بیند و به دلیل سادگی، سرعت استنتاج آن در عمل بسیار بالاست [\lr{30}]؛ نسخه بهبودیافته آن با معماری سبک‌تر و دقت بالاتر در \lr{TPAMI} منتشر شد [\lr{31}]. \lr{RUAS} با ترکیب \lr{unrolling} بازتعدی \lr{Retinex} و جست‌وجوی معماری، قابلیت تعمیم را در چارچوب \lr{zero-shot} ارتقا داد [\lr{21}]. نقطه عطف بعدی، \lr{SCI} بود: چارچوب یادگیری روشنایی خودکالیبره که با اشتراک وزن و ماژول خودکالیبراسیون، با تنها سه لایه کانولوشن ۳$\times$۳، هم پایداری نوردهی را در صحنه‌های متنوع تضمین کرد و هم در کاربردهای پایین‌دستی (تشخیص چهره در تاریکی روی \lr{DarkFace} و قطعه‌بندی معنایی شبانه روی \lr{ACDC)} برتری عملی خود را نشان داد [\lr{32}]. این ارزیابی دوگانه (کیفیت + وظیفه پایین‌دستی) الگویی است که رساله حاضر نیز آن را اتخاذ می‌کند. \lr{PairLIE} راهی برای یادگیری تنها از نمونه‌های کم‌نورِ جفت‌شده، بدون تصویر مرجع با نور عادی، پیشنهاد داد [\lr{33}] و \lr{SMG} با مدلسازی و هدایت ساختاری، جزئیات بازسازی‌شده را غنی‌تر ساخت [\lr{34}].\par
\begin{figure}[H]\centering
\includegraphics[width=0.92\linewidth]{figures/fig-08.png}\\[4pt]
\end{figure}
شکل ۸: چارچوب \lr{Zero-DCE} (برگرفته از \lr{Fig.2} در مرجع [\lr{30}])\par
\begin{figure}[H]\centering
\includegraphics[width=0.92\linewidth]{figures/fig-09.png}\\[4pt]
\end{figure}
شکل ۹: مقایسه کیفیت/کارایی روش‌های \lr{unsupervised} (برگرفته از \lr{Fig.1} در مرجع [\lr{32}])\par
جدول ۳ همین استدلال را در سطح هزینه محاسباتی کمّی می‌کند: \lr{SCI} با تنها ۰/۰۰۰۳ میلیون پارامتر و ۰/۰۶۲ گیگافلاپس، مرتبه‌هایی سبک‌تر از حتی \lr{Zero-DCE} است؛ ویژگی‌ای که استقرار روی دستگاه‌های لب کاربر را عملی می‌سازد و در فصل روش‌شناسی به آن بازخواهیم گشت.\par
جدول ۳: هزینه محاسباتی روش‌های منتخب: اندازه مدل، گیگافلاپس و زمان استنتاج روی \lr{GPU} (اعداد از [\lr{32}])\par
\begin{table}[H]\centering
{\small
\begin{tabular}{|c|c|c|c|}
\hline
روش & اندازه مدل (\lr{M)} & \lr{FLOPs} (\lr{G)} & زمان استنتاج (\lr{s)} \\ \hline
\lr{RetinexNet} & \lr{0.84} & \lr{136.02} & \lr{0.119} \\ \hline
\lr{KinD} & \lr{8.54} & \lr{29.13} & \lr{0.181} \\ \hline
\lr{DRBN} & \lr{0.58} & \lr{37.79} & \lr{0.053} \\ \hline
\lr{EnlightenGAN} & \lr{8.64} & \lr{61.01} & \lr{0.0097} \\ \hline
\lr{Zero-DCE} & \lr{0.079} & \lr{5.21} & \lr{0.0042} \\ \hline
\lr{RUAS} & \lr{0.0014} & \lr{0.28} & \lr{0.0063} \\ \hline
\lr{SCI} & \lr{0.0003} & \lr{0.062} & \lr{0.0017} \\ \hline
\end{tabular}}
\end{table}
جریان پژوهشی این خط همچنان زنده است: \lr{Self-DACE++} با منحنی‌های تطبیقی و آموزش با ترتیب تصادفی، به استنتاج بلادرنگ روی دستگاه‌های محدود رسیده است [\lr{35}] و \lr{InterLight} با استخراج و عملیاتی‌سازی \lr{priors} روشنایی ذاتی و حافظه گیت‌شده با درخشندگی، نشان داد که اولویت‌دهی بازسازی نواحی تاریک همراه با حفظ وفاداری نواحی روشن، مسیر نویدبخشی است [\lr{36}]؛ بینشی که مستقیماً با ماهیت «نوردهی نامتوازن» مرتبط است و در طراحی ماژول گیت روشنایی (بخش ۴) به کار گرفته می‌شود.\par
\textbf{۱-۳-۳. سقف اطلاعاتی روش‌های بدون جفت‌داده}\par
با وجود مزیت تعمیم‌پذیری، خانواده \lr{zero-shot} یک محدودیت ذاتی دارد که می‌توان آن را «سقف اطلاعاتی» نامید: چون این روش‌ها تنها از تصویر ورودی و قیدهای هندسی-آماری بهره می‌گیرند، هیچ دانشی درباره «شکل طبیعی صحنه‌های دنیا» در اختیار ندارند. نتیجه در عمل، سه الگوی تکرارشونده است: (الف) تقویت نویز نهفته در نواحی بسیار تاریک [\lr{35}]، (ب) انحراف و اشباع رنگ در صحنه‌های با منابع نوری ترکیبی [\lr{7}]، و (ج) عدم توانایی در بازسازی جزئیاتی که در تصویر ورودی اساساً ثبت نشده‌اند. به بیان دیگر، این روش‌ها می‌توانند «روشنایی موجود را بازتوزیع کنند» اما نمی‌توانند «اطلاعات گم‌شده را بازسازی کنند». پر کردن این شکاف مستلزم \lr{prior}ی قوی و آموخته‌شده از توزیع تصاویر طبیعی است؛ دقیقاً همان چیزی که مدل‌های مولد بزرگ فراهم می‌کنند و موضوع زیربخش بعدی است.\par
\textbf{۱-۴. موج \lr{diffusion-prior}: \lr{prior} مولد طبیعت به‌عنوان «دانش رایگان»}\par
\textbf{۱-۴-۱. چرا مدل‌های انتشار؟}\par
مدل‌های انتشار احتمالاتی [\lr{37}] با آموختن توزیع داده‌های تصویری در فرایند حذف تدریجی نویز، یک \lr{prior} مولد قدرتمند برای تصاویر طبیعی فراهم می‌کنند و نسخه‌های متن‌به‌تصویرِ پیش‌آموزش‌دیده روی مجموعه‌داده‌های وب‌مقیاس، دامنه این \lr{prior} را به صحنه‌های بسیار متنوع گسترش داده‌اند [\lr{29}]. برای \lr{LLIE} این به معنای دسترسی به \lr{prior}ی است که دو ضعف خطوط قبلی را هم‌زمان رفع می‌کند: برخلاف خط نظارت‌شده، نیاز به آموزش اختصاصی جفت‌محور را می‌توان حذف یا کاهش داد؛ و برخلاف سقف اطلاعاتی \lr{zero-shot}، دانش غنی از ظاهر طبیعی صحنه‌ها در اختیار است. مقاله مروری ۲۰۲۵ این حوزه، شش رده شامل تجزیه ذاتی، طیفی-نهفته، شتاب‌یافته، هدایت‌شده، چندوجهی و خودمختار را معرفی و چالش‌های استقرار (حافظه، تأخیر، انرژی) را تحلیل کرده است [\lr{38}].\par
\begin{figure}[H]\centering
\includegraphics[width=0.92\linewidth]{figures/fig-10.png}\\[4pt]
\end{figure}
شکل ۱۰: رده‌بندی شش‌گانه \lr{LLIE} مبتنی بر \lr{diffusion} (برگرفته از [\lr{38}])\par
\textbf{۱-۴-۲. نسل اول: \lr{diffusion} نظارت‌شده برای \lr{LLIE}}\par
\lr{Diff-Retinex} نخستین تلاش جدی برای ترکیب تئوری \lr{Retinex} و مدل مولد بود: به‌جای بهبود مستقیم، تجزیه و شتاب‌دهی جداگانه با شبکه‌های انتشار انجام می‌شود و با حفظ شبکه همساز، وفاداری محتوایی تأمین می‌شود [\lr{39}]. \lr{DiffLL} با انتقال فرایند انتشار به حوزه فرکانس پایین موجک، سرعت استنتاج را به‌طور چشمگیر بهبود داد [\lr{40}] و \lr{PyDiff} با هرم تدریجی رزولوشن و تصحیح‌گر سراسری رنگ، همین مسیر را ادامه داد [\lr{41}]. \lr{LLDiffusion} بازنمایی‌های اغتشاش را در فرایند انتشار یاد می‌گیرد [\lr{42}]، \lr{CLE-Diffusion} کنترل‌پذیری نوردهی را بهبود می‌بخشد [\lr{43}] و \lr{ExposureDiffusion} اصول نوردهی را مستقیماً در آموزش مدل انتشار تزریق کرد [\lr{44}]. بهبود کنترل شرطی‌سازی با \lr{ReCo-Diff} به اوج رسید: به‌جای تغذیه مستقیم نقشه‌های خام \lr{Retinex}، یک استراتژی شرط دو-مرحله‌ای (یادگیری شرط قابل‌اعتماد و پالایش ویژگی-و-تصویرمحور) این مدل را به \lr{PSNR} حدود ۲۷/۶ در \lr{LOL} و ۲۹/۳ در \lr{LOL-v2-real} رساند؛ بالاتر از تمامی روش‌های سابق [\lr{45}].\par
در همین راستا، \lr{GPP-LLIE} با استخراج \lr{priors} ادراکی مولد از مدل‌های زبان-بینایی و تزریق آن‌ها به لایه‌ی نرمال‌سازی و سازوکار توجهِ فرایند انتشار، تعمیم‌پذیری روی داده‌های واقعی را به‌طور معناداری بهبود داد [\lr{46}].\par
\begin{figure}[H]\centering
\includegraphics[width=0.92\linewidth]{figures/fig-11.png}\\[4pt]
\end{figure}
شکل ۱۱: مقایسه استراتژی‌های شرطی‌سازی در مدل‌های انتشار \lr{LLIE} (برگرفته از \lr{Fig.1} در مرجع [\lr{45}])\par
جدول ۴ سیر صعودی کمی این موج و فاصله‌اش با خطوط قبلی را روی \lr{LOL} و \lr{LOL-v2-real} خلاصه می‌کند؛ توجه شود که \lr{LPIPS} و کیفیت ادراکی، جایی که مدل‌های انتشار بیشترین برتری را دارند، محور مقایسه است.\par
جدول ۴: مقایسه خطوط \lr{unsupervised/zero-shot} و \lr{diffusion} روی \lr{LOL} و \lr{LOL-v2-real} (اعداد از گزارش \lr{ReCo-Diff} [\lr{45}])\par
\begin{table}[H]\centering
{\small
\begin{tabular}{|c|c|c|c|c|c|}
\hline
روش & \lr{LOL} \lr{PSNR} & \lr{LOL} \lr{SSIM} & \lr{LOL} \lr{LPIPS} & \lr{v2-real} \lr{PSNR} & \lr{v2-real} \lr{SSIM} \\ \hline
\lr{Zero-DCE} & \lr{14.861} & \lr{0.562} & \lr{0.335} & \lr{18.059} & \lr{0.580} \\ \hline
\lr{KinD} & \lr{20.870} & \lr{0.799} & \lr{0.207} & \lr{17.544} & \lr{0.669} \\ \hline
\lr{PairLIE} & \lr{19.510} & \lr{0.736} & \lr{0.248} & \lr{20.357} & \lr{0.782} \\ \hline
\lr{SMG} & \lr{23.684} & \lr{0.826} & \lr{0.118} & \lr{24.620} & \lr{0.867} \\ \hline
\lr{SNR-Net} & \lr{24.608} & \lr{0.840} & \lr{0.151} & \lr{21.479} & \lr{0.848} \\ \hline
\lr{LLFlow} & \lr{24.999} & \lr{0.870} & \lr{0.117} & \lr{26.200} & \lr{0.888} \\ \hline
\lr{Retinexformer} & \lr{25.153} & \lr{0.843} & \lr{0.131} & \lr{22.794} & \lr{0.839} \\ \hline
\lr{Diff-Retinex} & \lr{21.981} & \lr{0.863} & \lr{0.048} & - & - \\ \hline
\lr{DiffLL} & \lr{26.336} & \lr{0.845} & \lr{0.217} & \lr{28.857} & \lr{0.876} \\ \hline
\lr{PyDiff} & \lr{27.088} & \lr{0.875} & \lr{0.111} & \lr{27.236} & \lr{0.869} \\ \hline
\lr{ReCo-Diff} & \lr{27.626} & \lr{0.884} & \lr{0.090} & \lr{29.306} & \lr{0.906} \\ \hline
\end{tabular}}
\end{table}
\begin{figure}[H]\centering
\includegraphics[width=0.92\linewidth]{figures/fig-12.png}\\[4pt]
\end{figure}
شکل ۱۲: معماری \lr{ReCo-Diff} (برگرفته از \lr{Fig.2} در مرجع [\lr{45}])\par
\textbf{۱-۴-۳. نسل دوم: \lr{prior} انتشار بدون آموزش؛ زاویه \lr{zero-shot}}\par
نقطه عطف برای رساله حاضر، ظهور رویکردی است که از مدل انتشار پیش‌آموزش‌دیده صرفاً به‌عنوان «\lr{prior}» استفاده می‌کند و هیچ آموزش، بهینه‌سازی یا \lr{fine-tuning} نیاز ندارد. \lr{Cho} و همکاران نشان دادند که با هدایت استنتاج یک مدل متن‌به‌تصویر پیش‌آموزش‌دیده با ویژگی‌های درونی خود مدل، می‌توان تصاویر کم‌نور را با وفاداری بالا بازسازی کرد، بدون هیچ گام آموزشی، و حتی همین روش، بدون تغییر، در وظیفه تعادل سفیدی خودکار عملکرد نزدیک به \lr{SOTA} ارائه داد [\lr{29}]. \lr{He} و همکاران این ایده را با ادغام \lr{priors} حوزه فرکانس (موجک + فوریه) در فرایند نمونه‌گیری غنی‌تر کردند تا نقص راهنمایی نوری و ساختاری در سناریوهای با اغتشاش شدید جبران شود [\lr{47}]. در موازات، \lr{DarkDiff} نشان داد مدل انتشار پیش‌آموزش‌دیده را می‌توان برای \lr{ISP} دوربین «بازمالشافت» کرد و کیفیت ادراکی بهتری نسبت به مدل‌های رگرسیونی در سه بنچمارک \lr{RAW} کم‌نور به دست آورد [\lr{6}]، و \lr{OWDiff} با موجک هم‌پوشان، آرتیفکت‌های بلوکی نسل قبلی را ساختاریافته رفع کرد [\lr{48}].\par
\begin{figure}[H]\centering
\includegraphics[width=0.92\linewidth]{figures/fig-13.png}\\[4pt]
\end{figure}
شکل ۱۳: خط لوله \lr{DarkDiff} (برگرفته از \lr{Fig.2} در مرجع [\lr{6}])\par
این نسل دوم از سه منظر برای پایان‌نامه حاضر ایده‌آل است: (۱) نیاز محاسباتی آموزش را حذف می‌کند و تنها به استنتاج (و گاه بهینه‌سازی زمان-استنتاج) روی یک \lr{GPU} مصرفی متکی است؛ (۲) \lr{prior} مولد، مستقل از رژیم نوری دیتاست آموزشی است، پس پتانسیل تعمیم به رژیم‌های دیده‌نشده دارد؛ (۳) قابلیت افزودن \lr{priors} کمکی (فرکانسی، \lr{Retinex}، روشنایی ذاتی) معماری‌ای باز برای نوآوری فراهم می‌کند.\par
\textbf{۱-۴-۴. با این حال: تمام آزمایش‌ها هنوز روی «شب مطلق» است}\par
جمع‌بندی سنجیده‌ی موج \lr{diffusion}: با وجود انقلاب ادراکی، تقریباً تمام روش‌های این خانواده (از \lr{Diff-Retinex} تا \lr{ReCo-Diff} و حتی روش‌های \lr{zero-shot} [\lr{29}, \lr{47}]) روی بنچمارک‌های \lr{LOL/LOL-v2} ارزیابی شده‌اند؛ بنچمارک‌هایی که توزیعشان «تاریک یکنواخت شبانه» است. دو مسئله حل‌نشده باقی می‌ماند: نخست، تأخیر استنتاج مدل‌های انتشار برای کاربردهای بلادرنگ [\lr{38}]؛ و دوم، و مهم‌تر از همه، هیچ‌یک با رژیم نوری واسطی مواجه نشده‌اند که روشنایی و تاریکی را هم‌زمان و نامتوازن در یک قاب دارد. این دومین مسئله، موضوع زیربخش پایانی است.\par
\textbf{۱-۵. شکاف پژوهشی: رژیم نوری فراموش‌شده گرگ‌ومیش}\par
\textbf{۱-۵-۱. گرگ‌ومیش: تعریف علمی و چرایی سختی}\par
گرگ‌ومیش (\lr{Twilight)} بازه زمانی میان غروب و تاریکی کامل است که در سه فاز استاندارد طبقه‌بندی می‌شود: گرگ‌ومیش مدنی (خورشید تا ۶ درجه زیر افق)، دریایی (۶ تا ۱۲ درجه) و نجومی (۱۲ تا ۱۸ درجه)؛ در تمام این فازها، روشنایی صحنه در طول چند ده دقیقه، به‌سرعت و به‌شکل غیرخطی کاهش می‌یابد. از منظر پردازش تصویر، این رژیم ترکیبی به‌مراتب دشوارتر از تاریکی یکنواخت ایجاد می‌کند: عدم توازن نوردهی ذاتی: آسمانِ هنوز روشن با جاده تاریک، سایه‌های کشیده و بلند، گلوهای چراغ خودرو در حال ظهور و روشنایی‌های مصنوعی شهری که به‌تدریج روشن می‌شوند. برخلاف شب مطلق که یک «کمبود یکنواخت نور» است و می‌توان آن را با تقویت سراسری جبران کرد، گرگ‌ومیش «توزیع متضاد نور در یک قاب» است؛ هر تقویت سراسری یا آسمان را می‌سوزاند یا جاده را تاریک می‌گذارد.\par
این پدیده اخیراً در ادبیات به‌رسمیت شناخته شده است: \lr{UBLLIE} نشان داد تصاویر \lr{backlit} و \lr{low-light} با «عدم توازن شدید نوردهی» مواجه‌اند و روش‌های موجود هر دو رژیم را هم‌زمان پوشش نمی‌دهند [\lr{49}] و \lr{InterLight} با اولویت‌دهی بازسازی نواحی تاریک در کنار حفظ وفاداری نواحی روشن، عملاً به همین معماری دوگانه اعتراف کرده است [\lr{36}]. شکل ۱۴ نمونه‌هایی از این عدم توازن را نمایش می‌دهد. این پدیده پیشینه‌ای داوری‌شده نیز دارد: \lr{Afifi} و همکاران [\lr{50}] در \lr{CVPR} \lr{2021} تصحیح هم‌زمان خطاهای نوردهی کم و زیاد را صورت‌بندی کرده و دیتاست ۲۴هزارتایی \lr{MSEC} را معرفی کردند، و \lr{GPP-LLIE} در \lr{AAAI} \lr{2025} صریحاً گزارش داد که روش‌های موجود هنگام مواجهه با تصاویر واقعیِ متفاوت از شرایط آموزش، دچار «نتایج نامتوازن و آرتیفکت بازنوردهی» می‌شوند [\lr{46}]. گرگ‌ومیش بارزترین مصداق طبیعی این رده در صحنه‌های بیرونی است؛ اما همان‌طور که در ادامه می‌آید، از بنچمارک‌های استاندارد \lr{LLIE} غایب بوده است.\par
\begin{figure}[H]\centering
\includegraphics[width=0.92\linewidth]{figures/fig-14.png}\\[4pt]
\end{figure}
شکل ۱۴: نمونه‌های عدم توازن نوردهی (\lr{backlit} + \lr{low-light)} (برگرفته از \lr{Fig.1} در مرجع [\lr{49}])\par
\textbf{۱-۵-۲. غیاب در دیتاست‌ها}\par
مرور دیتاست‌های استاندارد \lr{LLIE}، تصویری گویا از این غیاب ترسیم می‌کند. \lr{LOL} و \lr{LOL-v2} جفت‌های شب مطلق داخلی/شهری‌اند [\lr{15}, \lr{27}]؛ \lr{SID}، \lr{SMID} و \lr{SDSD} جفت‌های \lr{RAW/}ویدئویی شبانه‌اند [\lr{16}, \lr{51}]؛ \lr{SICE} صحنه‌های نوردهی مخلوط دارد اما نه در سناریوی رانندگی و نه با تمرکز بر رژیم گرگ‌ومیش [\lr{52}]؛ و \lr{R2RNet} با دیتاست \lr{LSRW}، باز هم بر «\lr{real-low}» شبانه متمرکز است [\lr{53}]. در سمت رانندگی، \lr{DarkDriving} (روز/شب هم‌تراز) [\lr{5}]، \lr{LENVIZ} (نوردهی پایین شبانه) [\lr{54}] و \lr{LoLI-Street} (خیابانی کم‌نور) [\lr{55}] همگی رژیم «شب» را هدف گرفته‌اند و \lr{Dark} \lr{Zurich} و \lr{NightCity} نیز همین وضع را دارند. \lr{ACDC} چهار وضعیت آب‌وهوایی/نوری دارد اما گرگ‌ومیش را به‌صورت مستقل بنچمارک نمی‌کند [\lr{4}]. شکل ۱۵ نمونه‌های این بنچمارک‌ها را کنار هم می‌گذارد. نتیجه این مرور: تا حد اطلاع نویسندگان، هیچ مجموعه‌داده یا بنچمارک استانداردی برای بهبود تصویر در رژیم گرگ‌ومیش (به‌ویژه در سناریوی رانندگی) همراه با پروتکل ارزیابی اختصاصی این رژیم شناسایی نشد. مجموعه‌هایی مانند \lr{ExDark} [\lr{3}] گرچه صحنه‌های گرگ‌ومیش را در میان شرایط خود دارند، برای بهبود تصویر در این رژیم ساخته نشده‌اند و پروتکل ارزیابی مستقلی برای آن ارائه نمی‌دهند.\par
\begin{figure}[H]\centering
\includegraphics[width=0.92\linewidth]{figures/fig-15`.png}\\[4pt]
\includegraphics[width=0.92\linewidth]{figures/fig-15a.png}\\[4pt]
\includegraphics[width=0.92\linewidth]{figures/fig-15b.png}\\[4pt]
\end{figure}
شکل ۱۵: نمونه‌های بنچمارک‌های موجود (همگی شب مطلق) (ترکیب‌شده از [\lr{54}]، [\lr{55}]، [\lr{5}])\par
\begin{center}\fbox{جای تصاویر گرگ‌ومیش فاز ۳ ـــ \lr{TwilightDrive}}\end{center}
شکل ۱۶: نمونه‌های گرگ‌ومیش ضبط‌شده در این پژوهش (فاز ۳، دیتاست پیشنهادی \lr{TwilightDrive)}\par
\textbf{۱-۵-۳. غیاب در روش‌ها و نتیجه‌گیری}\par
در سمت روش‌ها وضعیت به همین منوال است: تا حد اطلاع نویسندگان، در مرور انجام‌شده هیچ‌یک از روش‌های مرورشده در این فصل (از \lr{RetinexNet} تا \lr{Retinexformer}، از \lr{Zero-DCE} و \lr{SCI} تا \lr{Diff-Retinex} و \lr{ReCo-Diff} و روش‌های \lr{zero-shot} مبتنی بر \lr{diffusion-prior)} برای رژیم گرگ‌ومیش طراحی، آموزش یا ارزیابی نشده‌اند [\lr{38}]. از جمع سه زیربخش اخیر، شکاف پژوهشی رساله حاضر با دقت قابل‌بیان است:\par
بر مبنای این شکاف، رساله حاضر سه مشارکت را دنبال می‌کند: (۱) توسعه یک روش بهبود کم‌نور \lr{zero-shot} مبتنی بر \lr{diffusion-prior} که با \lr{priors} روشنایی-آگاه (الهام از [\lr{36}] و [\lr{47}]) برای نوردهی نامتوازن گرگ‌ومیش سازگار شود؛ (۲) ساخت و انتشار \lr{TwilightDrive}، نخستین بنچمارک \lr{unpaired} گرگ‌ومیش رانندگی با پروتکل ضبط و کیوریشن استاندارد و یک زیرمجموعه مرجع‌دار (\lr{bracketed)} محدود؛ و (۳) ارزیابی اثر بهبود گرگ‌ومیش بر وظیفه پایین‌دستی تشخیص شیء در صحنه‌های رانندگی. این سه مشارکت، ساختار فرضیه‌ها (بخش ۲)، روش‌شناسی (بخش ۴) و نوآوری (بخش ۵) این طرح را تعیین می‌کنند.\par
\textbf{۱-۶. اهداف پژوهش}\par
هدف کلی: ارائه یک روش بهبود تصویر کم‌نور \lr{zero-shot} مبتنی بر \lr{prior} مدل‌های انتشار، سازگار با رژیم نوری گرگ‌ومیش، همراه با نخستین بنچمارک ارزیابی این رژیم، برای ارتقای پایداری سامانه‌های بینایی ماشین در شرایط نوری واسط.\par
اهداف جزئی:\par
۱. تحلیل کمّی رفتار روش‌های \lr{SOTA} \lr{zero-shot} و \lr{diffusion-prior} در رژیم گرگ‌ومیش و شناسایی سازوکارهای خرابی آن‌ها (نخستین مطالعه از این دست تا حد اطلاع نویسندگان).\par
۲. طراحی معماری هدایت‌شده با \lr{priors} روشنایی-آگاه در فرایند نمونه‌گیری انتشار برای حفظ تعادل نوردهی نواحی روشن/تاریک.\par
۳. جمع‌آوری، کیوریشن و انتشار دیتاست \lr{TwilightDrive} (۱۰۰۰+ تصویر، ۳ فاز گرگ‌ومیش، ۲ سنسور، همراه با زیرمجموعه \lr{bracketed).}\par
۴. سنجش اثر بهبود گرگ‌ومیش بر دقت تشخیص شیء به‌عنوان سنجه کاربردی نهایی (مطالعه‌ای که تا حد اطلاع نویسندگان پیشینه‌ای ندارد).\par


\section{سوالات تحقیق/فرضیه‌ها}
سوالات تحقیق/فرضیه‌ها\par
سوال اصلی تحقیق: چگونه می‌توان با بهره‌گیری از \lr{prior} مدل‌های انتشار پیش‌آموزش‌دیده، بهبود تصویر کم‌نور بدون جفت‌داده را برای رژیم نوری نامتوازن گرگ‌ومیش سازگار ساخت؟\par
سوالات فرعی:\par
ــ س۱. روش‌های موجود \lr{zero-shot} و \lr{diffusion-prior} در رژیم گرگ‌ومیش دقیقاً با چه سازوکاری دچار افت کیفیت می‌شوند؟\par
ــ س۲. کدام \lr{priors} کمکی (روشنایی-آگاه، فرکانسی) بیشترین اثر را بر تعادل نوردهی دارند؟\par
ــ س۳. بهبود گرگ‌ومیش چه اندازه بر عملکرد وظیفه پایین‌دستی (تشخیص شیء) اثر می‌گذارد؟\par
فرضیه‌ها (کمی و قابل‌آزمون):\par
فرضیه ۱ (کیفیت در رژیم گرگ‌ومیش): ادغام \lr{guidance} روشنایی-آگاه در فرایند نمونه‌گیری یک مدل انتشار پیش‌آموزش‌دیده، میانگین معیارهای بدون مرجع (\lr{NIQE}، \lr{MUSIQ)} را روی دیتاست پیشنهادی \lr{TwilightDrive} و مجموعه‌های \lr{unpaired} (\lr{ExDark}، \lr{DarkFace}، \lr{DICM/NPE/MEF)} نسبت به قوی‌ترین روش‌های \lr{zero-shot} موجود [\lr{29}, \lr{47}] حداقل ۱۰\% بهبود می‌دهد.\par
سنجه سنجش: مقایسه جفتی معنادار (۵ اجرای مستقل، آزمون \lr{t} زوجی، و \lr{p} کمتر از ۰/۰۵).\par
فرضیه ۲ (حفظ تعادل نوردهی): سازوکار گیت روشنایی، نسبت پیکسل‌های کم‌نورده/بازنورده خروجی را نسبت به \lr{baseline} حداقل ۳۰\% کاهش داده و میانگین روشنایی خروجی را در بازه هدف ۰/۴ تا ۰/۶ در کانال \lr{Y} (آستانه پیشنهادی این پژوهش، متناظر با روشنایی میانی صحنه) نگه می‌دارد، بدون افت معنادار کیفیت روی شب مطلق (اختلاف \lr{PSNR} کمتر از ۰/۵ \lr{dB} روی زیرمجموعه \lr{bracketed} در مقایسه با \lr{baseline).}\par
سنجه سنجش: شاخص نسبت پیکسل‌های \lr{over-/under-exposed} و \lr{PSNR/SSIM} روی زیرمجموعه مرجع‌دار.\par
فرضیه ۳ (فیدلیتی در ارزیابی مرجع‌دار): روی زیرمجموعه مرجع‌دار \lr{TwilightDrive} (حدود ۵۰ صحنه با چندنوردهی)، روش پیشنهادی به \lr{PSNR/SSIM} قابل‌رقابت با روش‌های نظارت‌شده \lr{SOTA} (اختلاف کمتر از ۱ \lr{dB} [\lr{25}]) و \lr{LPIPS} بهتر می‌رسد، بدون استفاده از هیچ داده جفتی در آموزش.\par
سنجه سنجش: مقایسه با \lr{Retinexformer} [\lr{25}] و \lr{LLFlow} [\lr{24}] در حالت \lr{cross-domain} (آموزش‌شده روی \lr{LOL}، آزمون روی گرگ‌ومیش).\par
فرضیه ۴ (اثر پایین‌دستی): بهبود گرگ‌ومیش، \lr{mAP@0.5} تشخیص شیء (مدل پیش‌آموزش‌دیده روی تصاویر روز) را روی تصاویر گرگ‌ومیش نسبت به تصاویر خام و نسبت به بهبود با روش‌های مقایسه، حداقل ۵ واحد درصد مطلق (افزایش مطلق، نه نسبی) افزایش می‌دهد.\par
سنجه سنجش: ارزیابی روی زیرمجموعه برچسب‌خورده (\lr{annotated)} \lr{TwilightDrive} (حدود ۳۰۰ فریم) و بنچمارک‌های انوتیت موجود (\lr{LoLI-Street} [\lr{55}]).\par
نکته طراحی: هر فرضیه به یکی از مشارکت‌های رساله گره خورده است (فرضیه‌های ۱ و ۲ به مشارکت ۱، فرضیه ۳ به مشارکت ۲، فرضیه ۴ به مشارکت ۳) و ابطال هر یک، مسیر اصلاح مشخصی در فازهای اجرا باز می‌گذارد.\par


\section{کاربردهای تحقیق و استفاده‌کنندگان از نتایج رساله}
کاربردهای تحقیق و استفاده‌کنندگان از نتایج\par
ــ سامانه‌های رانندگی خودران و \lr{ADAS}: بازه گرگ‌ومیش یکی از دوره‌های چالش‌برانگیز برای ادراک بصری در رانندگی است؛ بهبود پایدار دید در این رژیم می‌تواند ادراک بصری سامانه‌ها در شرایط واقعی را تقویت کند.\par
ــ نظارت تصویری شهری و صنعتی: بازسازی رخدادها در تصاویر ثبت‌شده هنگام غروب و سپیده‌دم.\par
ــ دوربین‌های خودرویی (دش‌کم) و تلفن همراه: بهبود بلادرنگ‌پذیر با \lr{GPU} مصرفی (بدون آموزش اختصاصی).\par
ــ استفاده‌کنندگان: پژوهشگران بینایی ماشین و \lr{low-level} \lr{vision} (از طریق انتشار دیتاست \lr{TwilightDrive} و کد)، توسعه‌دهندگان \lr{ADAS} و خودروسازان، مراکز کنترل ترافیک، و صنعت دوربین موبایل.\par


\section{روش شناسی (روش و ابزار جمع‌آوری، انجام و تحلیل اطلاعات تحقیق)}
روش شناسی (روش، ابزار و اصول اجرا، انجاز و تحلیل اطلاعات تحقیق)\par
\textbf{۴-۱. داده‌های پژوهش}\par
الف) بنچمارک‌های استاندارد (برای مقایسه‌پذیری با ادبیات):\par
ــ جفت‌دار: \lr{LOL-v1/v2-real/synthetic} [\lr{15}, \lr{27}]، \lr{LSRW} [\lr{53}]، \lr{FiveK} [\lr{56}]\par
ــ چندنوردهی: \lr{SICE} [\lr{52}]\par
ــ بدون مرجع: \lr{DICM}، \lr{NPE}، \lr{MEF}، \lr{VV}، \lr{ExDark} [\lr{3}]، \lr{DarkFace} [\lr{2}]\par
ب) دیتاست پیشنهادی \lr{TwilightDrive} (مشارکت دوم):\par
ــ پروتکل ضبط: سه فاز گرگ‌ومیش (مدنی، دریایی و نجومی)، شهر و بزرگراه، شرایط آب‌وهوایی متنوع، و دو سنسور (دوربین موبایل و دش‌کم خودرو)؛ ضبط از خودروی ساکن برای حذف ابهام حرکتی.\par
ــ هدف: ۱,۰۰۰ تا ۲,۰۰۰ فریم کیوریشده (بیش از دو برابر \lr{LOL-v2-real)} با ثبت \lr{EXIF} (\lr{ISO}، نوردهی، زمان).\par
ــ کیوریشن (\lr{curation)}: غربال فریم‌های تار، حذف تکراری‌ها، و ناشناس‌سازی خودکار پلاک و چهره (\lr{anonymization)} (تشخیص و محو) پیش از انتشار عمومی.\par
ــ زیرمجموعه مرجع‌دار (\lr{bracketed)}: حدود ۵۰ صحنه با \lr{bracketing} چندنوردهی (\lr{AEB)} برای محاسبه \lr{PSNR/SSIM.}\par
ــ زیرمجموعه برچسب‌خورده: حدود ۳۰۰ فریم برای ارزیابی تشخیص شیء.\par
\textbf{۴-۲. روش پیشنهادی (معماری کلی)}\par
پایه روش، یک مدل انتشار متن‌به‌تصویر پیش‌آموزش‌دیده است که وزن‌های آن ثابت نگه داشته می‌شود [\lr{29}] و سه ماژول نرم پیشنهادی روی فرایند نمونه‌گیری آن سوار می‌شود:\par
\lr{1.} ماژول گیت روشنایی: تخمین نقشه روشنایی/سایه ورودی و تولید \lr{guidance} مکانی برای تعادل نواحی روشن/تاریک (الهام از [\lr{36}])؛ پاسخ مستقیم به فرضیه ۲.\par
\lr{2.} ماژول \lr{prior} فرکانسی: هدایت در حوزه موجک برای حفظ جزئیات و مهار انحراف رنگ [\lr{47}, \lr{40}].\par
\lr{3.} ماژول استنتاج تطبیقی: تعداد گام‌های نمونه‌گیری بر اساس سطح تاریکی صحنه تنظیم می‌شود تا تأخیر کاربردی کنترل شود.\par
ارزیابی مطابق الگوی دومسیره \lr{SCI} انجام می‌شود: کیفیت ادراکی + وظیفه پایین‌دستی [\lr{32}].\par
\textbf{۴-۳. معیارهای ارزیابی و مدل‌های مقایسه}\par
ــ مرجع‌دار: \lr{PSNR}، \lr{SSIM} [\lr{57}]، \lr{LPIPS} [\lr{58}]\par
ــ بدون مرجع: \lr{NIQE} [\lr{59}]، \lr{MUSIQ} [\lr{60}]، نسبت پیکسل‌های \lr{over/under-exposed}، \lr{EME}\par
ــ پایین‌دستی: \lr{mAP@0.5} تشخیص شیء با مدل \lr{YOLO} پیش‌آموزش‌دیده روی روز\par
ــ مدل‌های مقایسه (۴ خانواده):\par
ــ نظارت‌شده \lr{SOTA}: \lr{Retinexformer} [\lr{25}]، \lr{LLFlow} [\lr{24}]\par
ــ \lr{unsupervised}: \lr{Zero-DCE} [\lr{30}]، \lr{SCI} [\lr{32}]، \lr{Self-DACE++} [\lr{35}]\par
ــ \lr{diffusion} نظارت‌شده: \lr{Diff-Retinex} [\lr{39}]، \lr{ReCo-Diff} [\lr{45}]\par
ــ \lr{zero-shot} \lr{prior}: \lr{Cho} \lr{et} \lr{al.} [\lr{29}]، \lr{He} \lr{et} \lr{al.} [\lr{47}]\par
ــ پروتکل آماری: هر آزمایش ۵ بار با \lr{seed} متفاوت؛ گزارش میانگین و انحراف معیار؛ آزمون معناداری زوجی؛ \lr{ablation} کامل برای هر سه ماژول.\par
\textbf{۴-۴. ابزار و سخت‌افزار}\par
پیاده‌سازی با \lr{Python/PyTorch} روی مدل انتشار متن‌باز پیش‌آموزش‌دیده؛ سخت‌افزار: \lr{GPU} مصرفی \lr{RTX} \lr{5060} و در آزمایش‌های سنگین‌تر، \lr{Colab} \lr{Pro} (\lr{L4)}؛ خط لوله کیوریشن دیتاست (تحلیل \lr{EXIF}، تشخیص تاری و ناشناس‌سازی) نیز با همین بستر توسعه می‌یابد.\par
\textbf{۴-۵. ریسک‌های اجرا و راهکار}\par


\section{جنبه نوآوری تحقیق}
جنبه نوآوری تحقیق\par
نوآوری این پژوهش در سه لایه متمایز از کارهای موجود قرار دارد:\par
لایه ۱ (روش): سازگارسازی \lr{diffusion-prior} برای رژیم نامتوازن. روش‌های \lr{zero-shot} مبتنی بر \lr{diffusion-prior} مانند [\lr{29}] و [\lr{47}] فرض کرده‌اند که اغتشاش ورودی «تاریک یکنواخت» است و سازوکاری برای مواجهه با تعارض نوردهیِ هم‌زمان روشن/تاریک در یک قاب ندارند. ماژول گیت روشنایی پیشنهادی، سازوکاری است که تا حد اطلاع نویسندگان پیشینه‌ای مشابه در این خانواده ندارد: هدایت مکان‌آگاه استنتاج مدل انتشار پیش‌آموزش‌دیده، برای آنکه نواحی روشن سوزانده و نواحی تاریک رها نشوند، بدون هیچ آموزش وزن.\par
لایه ۲ (داده): بنچمارک پیشنهادی گرگ‌ومیش. بنچمارک‌های رانندگی/کم‌نور موجود (\lr{LENVIZ} [\lr{54}]، \lr{LoLI-Street} [\lr{55}] و \lr{DarkDriving} [\lr{5}]) بر رژیم‌های تاریک/شبانه متمرکزند و هیچ‌کدام پروتکل ارزیابی مستقلی برای گرگ‌ومیش ارائه نمی‌دهند. \lr{TwilightDrive}، مجموعه‌داده پیشنهادی این رساله، طراحی شده است تا نخستین بنچمارک عمومیِ اختصاصی گرگ‌ومیش \lr{LLIE} باشد: سه فاز استاندارد گرگ‌ومیش با دو سنسور و زیرمجموعه ارجاع‌دار و انوتیت‌شده؛ نوعی زیرساخت ارزیابی که در صورت انتشار، برای کل جامعه پژوهشی \lr{LLIE} قابل استفاده خواهد بود.\par
لایه ۳ (ارزیابی): سنجش کاربردی در رژیم جدید. الگوی ارزیابی دومسیره \lr{SCI} [\lr{32}] (کیفیت ادراکی و وظیفه پایین‌دستی) را به رژیم گرگ‌ومیش می‌بریم (اجرایی که تا حد اطلاع نویسندگان پیشینه‌ای ندارد) تا سهم بهبود گرگ‌ومیش در زنجیره ادراک رانندگی کمّی شود.\par
جدول مقایسه با نزدیک‌ترین کارها:\par
\begin{table}[H]\centering
{\small
\begin{tabular}{|c|c|c|c|c|}
\hline
قابلیت & \lr{Cho} \lr{et} \lr{al.} [\lr{29}] & \lr{He} \lr{et} \lr{al.} [\lr{47}] & \lr{Self-DACE++} [\lr{35}] & روش پیشنهادی \\ \hline
بدون جفت‌داده & بله & بله & بله & بله \\ \hline
بدون آموزش وزن & بله & بله & خیر & بله \\ \hline
\lr{prior} فرکانسی & خیر & بله & خیر & بله \\ \hline
سازگاری با نوردهی نامتوازن & خیر & خیر & خیر & بله (گیت روشنایی) \\ \hline
ارزیابی در رژیم گرگ‌ومیش & خیر & خیر & خیر & بله \\ \hline
دیتاست اختصاصی رژیم واسط & خیر & خیر & خیر & بله (\lr{TwilightDrive)} \\ \hline
ارزیابی وظیفه پایین‌دستی در گرگ‌ومیش & خیر & خیر & خیر & بله \\ \hline
استنتاج تطبیقی سبک‌وزن & خیر & خیر & بله & بله \\ \hline
\end{tabular}}
\end{table}


\section{جدول زمان‌بندی انجام پروژه}
\textbf{برنامه زمان‌بندی}\par
زمان‌بندی اجرایی پروژه در قالب پنج فاز طی ۲۴ ماه مطابق جدول زیر است.\par
جدول ۵: برنامه زمان‌بندی پنج‌فازی اجرای پروژه (تألیف پژوهش حاضر)\par
\begin{table}[H]\centering
{\small
\begin{tabular}{|c|c|c|c|}
\hline
فاز & ماه & فعالیت & خروجی مورد انتظار \\ \hline
۰ & ۱ تا ۳ & تکمیل مرور ادبیات؛ پیاده‌سازی و بازتولید \lr{baseline}ها (\lr{Retinexformer}، \lr{SCI}، \lr{Cho} \lr{et} \lr{al.)}؛ ساخت مجموعه ارزیابی اولیه گرگ‌ومیش برای تحلیل خرابی & پروپوزال نهایی، کد \lr{baseline}، گزارش تحلیل خرابی (پاسخ س۱) \\ \hline
۱ & ۴ تا ۸ & طراحی و پیاده‌سازی ماژول گیت روشنایی و \lr{prior} فرکانسی؛ ارزیابی اولیه روی بنچمارک‌های موجود & مدل ساده، مقاله کنفرانسی \\ \hline
۲ & ۹ تا ۱۴ & توسعه کامل سه ماژول؛ \lr{ablation} و بهینه‌سازی استنتاج تطبیقی؛ ارزیابی جامع روی بنچمارک‌های استاندارد & نسخه میانی روش، پیش‌نویس مقاله مجله اول \\ \hline
۳ & ۱۵ تا ۱۹ & ضبط و کیوریشن \lr{TwilightDrive} (پروتکل بخش ۴-۱)؛ ارزیابی پایین‌دستی (تشخیص شیء) & مقاله مجله اول، نسخه اولیه دیتاست \\ \hline
۴ & ۲۰ تا ۲۴ & تکمیل، ناشناس‌سازی و در صورت احراز شرایط، انتشار دیتاست و کد؛ مقاله دوم؛ نگارش رساله & مقاله دوم، دیتاست منتشرشده و رساله تکمیل‌شده \\ \hline
\end{tabular}}
\end{table}
نقاط کنترل: پایان فاز ۱ = تصمیم \lr{go/no-go} بر اساس نتایج فرضیه ۱؛ پایان فاز ۳ = ارزیابی فرضیه‌های ۲ و ۳؛ پایان فاز ۴ = دفاع بر اساس دو مقاله و دیتاست منتشرشده.\par


\section{مراجع}
\begin{thebibliography}{60}
\bibitem{zheng2024comprehensive} \lr{S. Zheng, Y. Ma, J. Pan, C. Lu, and G. Gupta, "Low-Light Image and Video Enhancement: A Comprehensive Survey and Beyond," arXiv preprint arXiv:2212.10772, 2024.}
\bibitem{yang2020darkface} \lr{W. Yang et al., "Advancing Image Understanding in Poor Visibility Environments: A Collective Benchmark Study," IEEE Transactions on Image Processing, vol. 29, pp. 5737-5752, 2020.}
\bibitem{loh2019exdark} \lr{Y. P. Loh, and C. S. Chan, "Getting to Know Low-Light Images with the Exclusively Dark Dataset," Computer Vision and Image Understanding, vol. 178, pp. 30-42, 2019.}
\bibitem{sakaridis2021acdc} \lr{C. Sakaridis, D. Dai, and L. Van Gool, "ACDC: The Adverse Conditions Dataset with Correspondences for Semantic Driving Scene Understanding," in IEEE/CVF International Conference on Computer Vision (ICCV), 2021.}
\bibitem{wang2026darkdriving} \lr{W. Wang et al., "DarkDriving: A Real-World Day and Night Aligned Dataset for Autonomous Driving in the Dark Environment," in IEEE International Conference on Robotics and Automation (ICRA), 2026.}
\bibitem{zheng2025darkdiff} \lr{A. Y. Zheng, Y. Zhang, J. Hu, R. A. Yeh, and C. Chen, "DarkDiff: Advancing Low-Light Raw Enhancement by Retasking Diffusion Models for Camera ISP," arXiv preprint arXiv:2505.23743, 2025.}
\bibitem{li2021survey} \lr{C. Li et al., "Low-Light Image and Video Enhancement Using Deep Learning: A Survey," IEEE Transactions on Pattern Analysis and Machine Intelligence, vol. 44, no. 12, pp. 9396-9416, 2021.}
\bibitem{land1977retinex} \lr{E. H. Land, "The Retinex Theory of Color Vision," Scientific American, vol. 237, no. 6, pp. 108-128, 1977.}
\bibitem{pizer1987clahe} \lr{S. M. Pizer et al., "Adaptive Histogram Equalization and Its Variations," Computer Vision, Graphics, and Image Processing, vol. 39, no. 3, pp. 355-368, 1987.}
\bibitem{jobson1997msrcr} \lr{D. J. Jobson, Z. Rahman, and G. A. Woodell, "A Multiscale Retinex for Bridging the Gap Between Color Images and the Human Observation of Scenes," IEEE Transactions on Image Processing, vol. 6, no. 7, pp. 965-976, 1997.}
\bibitem{dong2011dehaze} \lr{X. Dong et al., "Fast Efficient Algorithm for Enhancement of Low Lighting Video," in IEEE International Conference on Multimedia and Expo (ICME), 2011.}
\bibitem{fu2016srie} \lr{X. Fu, D. Zeng, Y. Huang, X. Zhang, and X. Ding, "A Weighted Variational Model for Simultaneous Reflectance and Illumination Estimation," in IEEE Conference on Computer Vision and Pattern Recognition (CVPR), 2016.}
\bibitem{guo2017lime} \lr{X. Guo, Y. Li, and H. Ling, "LIME: Low-Light Image Enhancement via Illumination Map Estimation," IEEE Transactions on Image Processing, vol. 26, no. 2, pp. 982-993, 2017.}
\bibitem{lore2017llnet} \lr{K. G. Lore, A. Akintayo, and S. Sarkar, "LLNet: A Deep Autoencoder Approach to Natural Low-Light Image Enhancement," Pattern Recognition, vol. 61, pp. 650-662, 2017.}
\bibitem{wei2018retinexnet} \lr{C. Wei, W. Wang, W. Yang, and J. Liu, "Deep Retinex Decomposition for Low-Light Enhancement," in British Machine Vision Conference (BMVC), 2018.}
\bibitem{chen2018sid} \lr{C. Chen, Q. Chen, J. Xu, and V. Koltun, "Learning to See in the Dark," in IEEE/CVF Conference on Computer Vision and Pattern Recognition (CVPR), 2018.}
\bibitem{zhang2019kind} \lr{Y. Zhang, J. Zhang, and X. Guo, "Kindling the Darkness: A Practical Low-Light Image Enhancer," in ACM International Conference on Multimedia (ACM MM), 2019.}
\bibitem{zhang2021beyond} \lr{Y. Zhang, X. Guo, J. Ma, W. Liu, and J. Zhang, "Beyond Brightening Low-Light Images," International Journal of Computer Vision, vol. 129, no. 4, pp. 1013-1037, 2021.}
\bibitem{wang2019deepupe} \lr{R. Wang, Q. Zhang, C. Fu, X. Shen, W. Zheng, and J. Jia, "Underexposed Photo Enhancement Using Deep Illumination Estimation," in IEEE/CVF Conference on Computer Vision and Pattern Recognition (CVPR), 2019.}
\bibitem{yang2020drbn} \lr{W. Yang, S. Wang, Y. Fang, Y. Wang, and J. Liu, "From Fidelity to Perceptual Quality: A Semi-Supervised Approach for Low-Light Image Enhancement," in IEEE/CVF Conference on Computer Vision and Pattern Recognition (CVPR), 2020.}
\bibitem{liu2021ruas} \lr{R. Liu, L. Ma, J. Zhang, X. Fan, and Z. Luo, "Retinex-Inspired Unrolling with Cooperative Prior Architecture Search for Low-Light Image Enhancement," in IEEE/CVF Conference on Computer Vision and Pattern Recognition (CVPR), 2021.}
\bibitem{wu2022uretinex} \lr{W. Wu, J. Weng, P. Zhang, X. Wang, W. Yang, and J. Jiang, "URetinex-Net: Retinex-Based Deep Unfolding Network for Low-Light Image Enhancement," in IEEE/CVF Conference on Computer Vision and Pattern Recognition (CVPR), 2022.}
\bibitem{xu2022snr} \lr{X. Xu, R. Wang, C. Fu, and J. Jia, "SNR-Aware Low-Light Image Enhancement," in IEEE/CVF Conference on Computer Vision and Pattern Recognition (CVPR), 2022.}
\bibitem{wang2022llflow} \lr{Y. Wang, R. Wan, W. Yang, H. Li, L. Chau, and A. C. Kot, "Low-Light Image Enhancement with Normalizing Flow," in AAAI Conference on Artificial Intelligence (AAAI), 2022.}
\bibitem{cai2023retinexformer} \lr{Y. Cai, H. Bian, J. Lin, H. Wang, R. Timofte, and Y. Zhang, "Retinexformer: One-Stage Retinex-Based Transformer for Low-Light Image Enhancement," in IEEE/CVF International Conference on Computer Vision (ICCV), 2023.}
\bibitem{yan2025hvi} \lr{Q. Yan et al., "HVI: A New Color Space for Low-Light Image Enhancement," in IEEE/CVF Conference on Computer Vision and Pattern Recognition (CVPR), 2025.}
\bibitem{yang2021lolv2} \lr{W. Yang, W. Wang, H. Huang, S. Wang, and J. Liu, "Sparse Gradient Regularized Deep Retinex Network for Robust Low-Light Image Enhancement," IEEE Transactions on Image Processing, vol. 30, pp. 2072-2086, 2021.}
\bibitem{jiang2021enlightengan} \lr{Y. Jiang et al., "EnlightenGAN: Deep Light Enhancement Without Paired Supervision," IEEE Transactions on Image Processing, vol. 30, pp. 2340-2349, 2021.}
\bibitem{cho2024zeroshot} \lr{J. Cho, S. Aghajanzadeh, Z. Zhu, and D. A. Forsyth, "Zero-Shot Low Light Image Enhancement with Diffusion Prior," arXiv preprint arXiv:2412.13401, 2024.}
\bibitem{guo2020zerodce} \lr{C. Guo et al., "Zero-Reference Deep Curve Estimation for Low-Light Image Enhancement," in IEEE/CVF Conference on Computer Vision and Pattern Recognition (CVPR), 2020.}
\bibitem{li2022zerodcepp} \lr{C. Li, C. Guo, and C. C. Loy, "Learning to Enhance Low-Light Image via Zero-Reference Deep Curve Estimation," IEEE Transactions on Pattern Analysis and Machine Intelligence, vol. 44, no. 8, pp. 4225-4238, 2022.}
\bibitem{ma2022sci} \lr{L. Ma, T. Ma, R. Liu, X. Fan, and Z. Luo, "Toward Fast, Flexible, and Robust Low-Light Image Enhancement," in IEEE/CVF Conference on Computer Vision and Pattern Recognition (CVPR), 2022.}
\bibitem{fu2023pairlie} \lr{Z. Fu, Y. Yang, X. Tu, Y. Huang, X. Ding, and K. Ma, "Learning a Simple Low-Light Image Enhancer From Paired Low-Light Instances," in IEEE/CVF Conference on Computer Vision and Pattern Recognition (CVPR), 2023.}
\bibitem{xu2023smg} \lr{X. Xu, R. Wang, and J. Lu, "Low-Light Image Enhancement via Structure Modeling and Guidance," in IEEE/CVF Conference on Computer Vision and Pattern Recognition (CVPR), 2023.}
\bibitem{wen2026selfdace} \lr{J. Wen et al., "Self-DACE++: Robust Low-Light Enhancement via Efficient Adaptive Curve Estimation," in arXiv preprint arXiv:2604.25367, 2026.}
\bibitem{wang2026interlight} \lr{Z. Wang, X. Zhang, L. Chang, S. Chen, J. Ma, and H. Zhang, "InterLight: Leveraging Intrinsic Illumination Priors for Low-Light Image Enhancement," in International Joint Conference on Artificial Intelligence (IJCAI), 2026.}
\bibitem{ho2020ddpm} \lr{J. Ho, A. Jain, and P. Abbeel, "Denoising Diffusion Probabilistic Models," in Advances in Neural Information Processing Systems (NeurIPS), 2020.}
\bibitem{adhikarla2025taxonomy} \lr{E. Adhikarla, Y. Liu, and B. D. Davison, "Diffusion Models for Low-Light Image Enhancement: A Multi-Perspective Taxonomy and Performance Analysis," arXiv preprint arXiv:2510.05976, 2025.}
\bibitem{yi2023diffretinex} \lr{X. Yi, H. Xu, H. Zhang, L. Tang, and J. Ma, "Diff-Retinex: Rethinking Low-Light Image Enhancement with a Generative Diffusion Model," in IEEE/CVF International Conference on Computer Vision (ICCV), 2023.}
\bibitem{jiang2023diffll} \lr{H. Jiang, A. Luo, S. Han, H. Fan, and S. Liu, "Low-Light Image Enhancement with Wavelet-Based Diffusion Models," in SIGGRAPH Asia, 2023.}
\bibitem{zhou2023pydiff} \lr{D. Zhou, Z. Yang, and Y. Yang, "Pyramid Diffusion Models for Low-Light Image Enhancement," in International Joint Conference on Artificial Intelligence (IJCAI), 2023.}
\bibitem{wang2023lldiffusion} \lr{T. Wang et al., "LLDiffusion: Learning Degradation Representations in Diffusion Models for Low-Light Image Enhancement," arXiv preprint arXiv:2307.14659, 2023.}
\bibitem{yin2023cle} \lr{Y. Yin, D. Xu, C. Tan, P. Liu, Y. Zhao, and Y. Wei, "CLE Diffusion: Controllable Light Enhancement Diffusion Model," in ACM International Conference on Multimedia (ACM MM), 2023.}
\bibitem{wang2023exposure} \lr{Y. Wang et al., "ExposureDiffusion: Learning to Expose for Low-Light Image Enhancement," in IEEE/CVF International Conference on Computer Vision (ICCV), 2023.}
\bibitem{wu2023recodiff} \lr{Y. Wu et al., "ReCo-Diff: Explore Retinex-Based Condition Strategy in Diffusion Model for Low-Light Image Enhancement," arXiv preprint arXiv:2312.12826, 2023.}
\bibitem{zhou2025gppllie} \lr{H. Zhou, W. Dong, X. Liu, Y. Zhang, G. Zhai, and J. Chen, "Low-Light Image Enhancement via Generative Perceptual Priors," in AAAI Conference on Artificial Intelligence (AAAI), 2025.}
\bibitem{he2024frequency} \lr{J. He, S. Palaiahnakote, A. Ning, and M. Xue, "Zero-Shot Low-Light Image Enhancement via Joint Frequency Domain Priors Guided Diffusion," arXiv preprint arXiv:2411.13961, 2024.}
\bibitem{peng2026owdiff} \lr{F. Peng, T. Suzuki, and S. Kyochi, "Overlapped Wavelet Diffusion for Low-Light Image Enhancement," IEICE Transactions on Information and Systems, 2026.}
\bibitem{yasin2026ubllie} \lr{Y. Yasin, M. Usman, I. Radwan, and S. Anwar, "UBLLIE: Unified Backlight and Low-Light Image Enhancement," arXiv preprint arXiv:2608.04429, 2026.}
\bibitem{afifi2021exposure} \lr{M. Afifi, K. G. Derpanis, B. Ommer, and M. S. Brown, "Learning Multi-Scale Photo Exposure Correction," in IEEE/CVF Conference on Computer Vision and Pattern Recognition (CVPR), 2021.}
\bibitem{wang2021sdsd} \lr{R. Wang, X. Xu, C. Fu, J. Lu, B. Yu, and J. Jia, "Seeing Dynamic Scene in the Dark: A High-Quality Video Dataset with Mechatronic Alignment," in IEEE/CVF International Conference on Computer Vision (ICCV), 2021.}
\bibitem{cai2018sice} \lr{J. Cai, S. Gu, and L. Zhang, "Learning a Deep Single Image Contrast Enhancer from Multi-Exposure Images," IEEE Transactions on Image Processing, vol. 27, no. 4, pp. 2049-2062, 2018.}
\bibitem{hai2023r2rnet} \lr{J. Hai et al., "R2RNet: Low-Light Image Enhancement via Real-Low to Real-Normal Network," Journal of Visual Communication and Image Representation, vol. 90, pp. 103712, 2023.}
\bibitem{aithal2025lenviz} \lr{M. Aithal et al., "LENVIZ: A High-Resolution Low-Exposure Night Vision Benchmark Dataset," arXiv preprint arXiv:2503.19804, 2025.}
\bibitem{islam2024lolistreet} \lr{M. T. Islam, I. Alam, S. S. Woo, S. Anwar, I. H. Lee, and K. Muhammad, "LoLI-Street: Benchmarking Low-Light Image Enhancement and Beyond," in Asian Conference on Computer Vision (ACCV), 2024.}
\bibitem{bychkovsky2011fivek} \lr{V. Bychkovsky, S. Paris, E. Chan, and F. Durand, "Learning Photographic Global Tonal Adjustment with a Database of Input/Output Image Pairs," in IEEE Conference on Computer Vision and Pattern Recognition (CVPR), 2011.}
\bibitem{wang2004ssim} \lr{Z. Wang, A. C. Bovik, H. R. Sheikh, and E. P. Simoncelli, "Image Quality Assessment: From Error Visibility to Structural Similarity," IEEE Transactions on Image Processing, vol. 13, no. 4, pp. 600-612, 2004.}
\bibitem{zhang2018lpips} \lr{R. Zhang, P. Isola, A. A. Efros, E. Shechtman, and O. Wang, "The Unreasonable Effectiveness of Deep Features as a Perceptual Metric," in IEEE/CVF Conference on Computer Vision and Pattern Recognition (CVPR), 2018.}
\bibitem{mittal2013niqe} \lr{A. Mittal, R. Soundararajan, and A. C. Bovik, "Making a ``Completely Blind'' Image Quality Analyzer," IEEE Signal Processing Letters, vol. 20, no. 3, pp. 209-212, 2013.}
\bibitem{ke2021musiq} \lr{J. Ke, Q. Wang, Y. Wang, P. Milanfar, and F. Yang, "MUSIQ: Multi-Scale Image Quality Transformer," in IEEE/CVF International Conference on Computer Vision (ICCV), 2021.}
\end{thebibliography}
\end{document}